\section{Resource-normalized measurement and the bi-factor law}
\label{sec:model}

Capacity planning benefits from a calibrated exponent that reflects consensus
cost rather than host contention or load-generator growth. We introduce an
isolation protocol, a bi-factor characterization, and an identifiability result
under stated assumptions.

\subsection{Resource-normalized measurement as an isolation tool}

If $\n$ validators share a host of $Q$ cores, each receives $Q/\n$ cores, so
measured throughput mixes algorithmic cost with a shrinking compute budget.
\RNM{} is designed to isolate those effects under controlled quotas.

\begin{Def}[Resource-normalized measurement]
\label{def:rnm}
A campaign follows the \emph{\RNM{} protocol} with quota~$\q$ if every validator
container receives the same CPU quota and memory limit independently of~$\n$,
with aggregate reservation $\n\q$ below host capacity. Normalized throughput is
$\TPSn\eqdef\TPS/\q$.
\end{Def}

\begin{Assumption}[Separable compute budget]
\label{as:separable}
On the operating range of quotas in the campaign,
$\TPS(\n,\cc,\q)=\q\cdot\Phi(\n,\cc)\cdot(1+r)$ with $|r|\le C\q$ for a
constant~$C$ independent of~$(\n,\cc)$ there.
\end{Assumption}

\paragraph{Assumption check.}
Under Assumption~\ref{as:separable} and a power-law factor
$\Phi\propto\n^{-\beta}$, the OLS slope of $\log\TPSn$ on $\log\n$ at fixed
concurrency is independent of~$\q$: the relative remainder is absorbed by the
intercept. Section~\ref{sec:exp} tests this invariance across
$\q\in\{0.20,0.25,0.33\}$. Within the evaluated deployment the calibrated model
is reported at the reference quota $\q=0.25$; outside that regime the numerical
coefficients must be re-fit.

\begin{Def}[Bi-factor scaling characterization]
\label{def:bifactor}
Under the \RNM{} protocol,
\begin{equation}
\TPS(\n,\cc)=\Tzero\cdot\gfun(\cc)\cdot\n^{-\beta},
\quad
\gfun(\cc)=\frac{\cc^{\gamma}}{1+(\cc/\csat)^{\theta}},
\label{eq:bifactor}
\end{equation}
with $\Tzero>0$, $\beta>0$, $\gamma\in(0,1]$, $\theta>\gamma$, $\csat>0$.
In the unsaturated regime $\cc\ll\csat$ this is log-linear:
$\log\TPS=\log\Tzero+\gamma\log\cc-\beta\log\n$.
\end{Def}

The form separates concurrency from validator-count scaling, reduces to a
power law for $\cc\ll\csat$, saturates at large~$\cc$, and contains the
classical single-factor fit as the case $\gamma=0$ (or $\lambda=0$ with
$\alpha=-\beta$).

\subsection{Identifiability}

Proposition~\ref{prop:identifiability} is analytical under the unsaturated
bi-factor model and stated path assumptions. Its importance is interpretive:
it shows how to read a measured single-factor exponent under a stated
concurrency path, rather than introducing a fundamentally new mathematical
identity. Section~\ref{sec:exp} checks consistency within the evaluated
deployment; that check does not extend the claim beyond the model assumptions.

\begin{Proposition}[Identifiability of the consensus exponent]
\label{prop:identifiability}
Let observations follow the unsaturated bi-factor law with additive noise of
mean zero and finite variance, along a concurrency path
$\log\cc_{i}=\lambda\log\n_{i}+\mu+\eta_{i}$
($\E[\eta_{i}\mid\n_{i}]=0$). If the empirical second moment of
$x_{i}=\log\n_{i}$ converges to a positive limit, the OLS slope $\ahat_{m}$ of
$\log\TPS_{i}$ on $x_{i}$ alone satisfies
$\ahat_{m}\xrightarrow{\Prob}\gamma\lambda-\beta$.
In particular $\ahat>0\iff\lambda>\beta/\gamma$, and $\ahat=-\beta\iff\lambda=0$.
\end{Proposition}

\begin{Proof}
Write $x_{i}=\log\n_{i}$, $y_{i}=\log\TPS_{i}$. The path reduces the unsaturated
law to $y_{i}=a+\alpha x_{i}+\varepsilon_{i}$ with $\E[\varepsilon_{i}\mid x_{i}]=0$,
finite variance, and
\begin{equation}
\alpha=-\beta+\gamma\lambda.
\label{eq:biasdecomp}
\end{equation}
Centred OLS gives
$\ahat_{m}=\alpha+N_{m}/D_{m}$ with
$N_{m}=\sum(x_{i}-\bar x)\varepsilon_{i}$ and
$D_{m}=\sum(x_{i}-\bar x)^{2}$.
The second-moment assumption yields $D_{m}/m\to\Var^{*}(x)>0$; the summands of
$N_{m}$ are uncorrelated, mean zero, and of bounded second moment, so
$N_{m}/m\to0$ in probability. Continuous mapping:
$\ahat_{m}\xrightarrow{\Prob}\alpha$. Signs follow from $\gamma>0$.
\end{Proof}

Thus $\alpha=\underbrace{-\beta}_{\text{protocol}}
+\underbrace{\gamma\lambda}_{\text{measurement}}$.
Only a factorial design ($\lambda=0$) recovers the consensus exponent under the
stated model. For illustration, $\gamma=0.8$, $\beta=0.5$ and $\lambda=1$ yield
$\ahat=+0.3$, and extrapolating from $\n=5$ to $\n=60$ differs by a factor
$\approx7$, matching the order of the discrepancy in~\cite{peniak2026a}.

\begin{Corollary}[Non-transferability of single-factor exponents]
\label{cor:sign}
Under the unsaturated bi-factor model, laboratories measuring along
$\lambda_{1}\ne\lambda_{2}$ obtain exponents differing by
$\gamma(\lambda_{1}-\lambda_{2})$. Without reporting~$\lambda$, published
single-factor exponents should not be treated as transferable protocol
properties across measurement designs.
\end{Corollary}

\subsection{Implementation}

On a factorial grid ($\lambda=0$) we estimate $(\log\Tzero,\gamma,\beta)$ by
OLS in log-space and $(\csat,\theta)$ by nonlinear least squares on the full
$\cc$-sweep (Section~\ref{sec:exp}). The calibrated model then feeds prediction
intervals and the sizing program of the next section within the range supported
by the campaign.
